\section{A common escape cone and integral descent}\label{sec:escape}

For $P=(x,y,z)$ write $\norm P=\max(|x|,|y|,|z|)$ and put
\begin{equation}\label{eq:domains}
 \begin{aligned}
 \Ssmall&=\{P\in\Z^3:\min(|x|,|y|,|z|)\le1\},\\
 \Cplus&=\{P\in\Z^3:|x|,|y|\ge2,\
 |z|\ge\max(|x|,|y|),\ |z|>2,\ xyz>0\}.
 \end{aligned}
\end{equation}
Both sets are invariant under the coordinate exchange $J$.
The gap between integer absolute values at most one and at least
two is used essentially in the following arguments.

\begin{lemma}[Common forward cone]\label{lem:cone}
Each letter maps $\Cplus$ into itself. Along any infinite
letter sequence in which both letters occur infinitely often,
a starting point in $\Cplus$ has third-coordinate modulus
tending to infinity.
\end{lemma}

\begin{proof}
Put $(a,b,c)=(|x|,|y|,|z|)$. Since $xyz>0$, the numbers $xz$
and $y$ have the same sign. Under $A$ the third modulus is
$ac-b>0$, and
\[
 ac-b\ge ac-c=(a-1)c\ge c.
\]
The first two new moduli are $a,c$, and their product with the
new third coordinate is positive. Thus the image lies in
$\Cplus$. Equality $ac-b=c$ requires and is equivalent to
$a=2$, $b=c$. Applying $J$ gives the $B$ statement, with
equality exactly when $b=2$, $a=c$.

The successive third moduli are nondecreasing integers. If
bounded, they eventually equal a fixed $c>2$. A late $A$ step
then forces the absolute triple $(2,c,c)$ and leaves it
unchanged. Any intervening $A$ steps do the same. The next
$B$ step gives third modulus $c^2-2>c$, a contradiction.
Such a later $B$ exists by hypothesis. Hence the moduli
tend to infinity.
\end{proof}

\begin{lemma}[All-large descent]\label{lem:descent}
A forward letter trajectory outside $\E$ that never enters
$\Cplus$ reaches $\Ssmall$ after finitely many steps.
\end{lemma}

\begin{proof}
While the trajectory is outside $\Ssmall$, all three moduli
$a,b,c$ are at least two. If $xyz<0$, the new third modulus
is $ac+b$ under $A$ and $bc+a$ under $B$; either image has
positive product and belongs to $\Cplus$. Thus this case
cannot persist on the stipulated trajectory.

Suppose $xyz>0$. If $c\ge\max(a,b)$, the point is in
$\Cplus$ unless $a=b=c=2$, in which case it belongs to
$\E$. Both alternatives are excluded. If $a=b>c$, either
letter gives third modulus $a(c-1)\ge a>2$ and enters the
cone.

It remains, after possibly applying $J$, that
$a>\max(b,c)$. An $A$ step gives
$ac-b>a(c-1)\ge a$ and enters the cone. For a $B$ step let
$d=bc-a$ be the signed new magnitude relative to the sign
of $x$. There are four exhaustive possibilities:
\begin{enumerate}
\item If $d\le-2$, the image has all moduli at least two
and negative product; the following letter enters the cone.
\item If $-1\le d\le1$, the image belongs to $\Ssmall$.
\item If $d\ge\max(b,c)$, the image belongs to the cone
unless $b=c=d=2$. That exceptional image would be in
$\E$, which has no preimage outside $\E$ by \eqref{eq:E}.
\item Otherwise $2\le d<\max(b,c)<a$, and the new maximum
is strictly smaller than the old maximum $a$.
\end{enumerate}
The strict second-coordinate maximum is exactly the conjugate
case under $J$, with $A$ and $B$ exchanged. Thus, as long as
the stipulated trajectory remains outside $\Ssmall$, the
only surviving case is strict descent of a positive integer
maximum. It cannot continue indefinitely.
\end{proof}

\begin{lemma}[Small-set exits]\label{lem:exit}
If $P\in\Ssmall$ and either letter takes $P$ outside
$\Ssmall$, its image belongs to $\Cplus$.
\end{lemma}

\begin{proof}
For $A$, such an exit requires $|x|,|z|\ge2$ and $|y|\le1$.
The product $xz$ strictly dominates $y$, so the image has
positive coordinate product and
\[
 |xz-y|\ge |x||z|-1>\max(|x|,|z|),\qquad |xz-y|\ge3.
\]
These are the cone inequalities. The $B$ case follows by $J$.
\end{proof}

\begin{corollary}\label{cor:periodicsmall}
Every letter phase in the expansion of a periodic admissible
word orbit lies either in $\E$ or in $\Ssmall$. In the latter
case no phase has the third coordinate as its only small
coordinate.
\end{corollary}

\begin{proof}
Expand the ordinary word cycle into all chronological letter
phases. This is a finite repeating trajectory whose schedule
contains both letters. Lemma~\ref{lem:cone} excludes the cone.
If one phase is in $\E$, all are, by invariance and
bijectivity. Otherwise Lemma~\ref{lem:descent} forces any
all-large phase to reach $\Ssmall$, which cannot subsequently
be left by Lemma~\ref{lem:exit}. A finite repeating trajectory
therefore has all its phases in $\Ssmall$.

If $|x|,|y|\ge2$ and $|z|\le1$, both possible predecessor
phases in \eqref{eq:inverses} are outside $\Ssmall$:
\[
 |xy-z|\ge |x||y|-1\ge3.
\]
Such a phase cannot occur. This directional predecessor
constraint is essential to the precise line count.
\end{proof}
